% !TeX root = ../main.tex
% ====================
% V. KET QUA THUC NGHIEM
% ====================
\section{Kết quả thực nghiệm}
\label{sec:experimental}

Bảng~\ref{tab:current_results} trình bày kết quả của bốn cấu hình trong cùng
thiết lập huấn luyện từ khởi tạo ImageNet trên Market-1501. Tất cả cấu hình
đều được huấn luyện $120$ epoch với batch size $64$ và learning rate $0.008$.

\begin{table*}[!t]
\centering
\renewcommand{\arraystretch}{1.12}
\setlength{\tabcolsep}{4pt}
\begin{tabular}{@{}L{4.0cm}ccccccc@{}}
\toprule
\textbf{Mô hình} & \textbf{mAP (\%)} & \textbf{Rank-1 (\%)} & \textbf{Rank-5 (\%)} & \textbf{Rank-10 (\%)} & \textbf{Độ trễ (ms)} & \textbf{FPS} & \textbf{\(\Delta\)mAP} \\
\midrule
Baseline TransReID & 89.2 & 95.1 & 98.3 & 98.9 & 4.626 & 216.1 & -- \\
Semantic alignment & 89.5 & 95.3 & 98.5 & 99.0 & 4.637 & 215.6 & +0.3 \\
Local reliability & 89.9 & 95.2 & \textbf{98.6} & \textbf{99.1} & 4.786 & 208.9 & +0.7 \\
Semantic alignment + local reliability & \textbf{90.1} & \textbf{95.4} & \textbf{98.6} & \textbf{99.1} & 4.783 & 209.1 & \textbf{+0.9} \\
\bottomrule
\end{tabular}
\caption{Kết quả so sánh công bằng trên Market-1501. \(\Delta\)mAP được tính
so với baseline TransReID trong cùng thiết lập.}
\label{tab:current_results}
\end{table*}

\subsection{Hiệu quả của semantic alignment}

Semantic alignment đạt $89.5\%$ mAP và $95.3\%$ Rank-1, cao hơn baseline lần
lượt $0.3$ và $0.2$ điểm phần trăm. Mức cải thiện này không lớn, nhưng có ý
nghĩa vì độ trễ gần như không thay đổi: thời gian suy luận tăng từ $4.626$ ms
lên $4.637$ ms, tương đương khoảng $0.23\%$. Kết quả này phù hợp với thiết kế
của phương pháp, vì semantic mask chỉ được dùng để tạo patch-level supervision
trong huấn luyện và không được dùng làm đầu vào ở pha suy luận.

Điểm quan trọng là semantic alignment không can thiệp trực tiếp vào global
branch. Phiên bản hiện tại chỉ dùng semantic như một regularizer ở không gian
patch token: patch token được khuyến khích giữ thông tin part-level trong quá
trình huấn luyện, nhưng descriptor cuối vẫn được tạo bởi pipeline TransReID/JPM
gốc. Cách tích hợp này giúp bổ sung tín hiệu căn chỉnh ngữ nghĩa mà không làm
thay đổi đường truy hồi chính.

\subsection{Hiệu quả của local reliability}

Local reliability đạt $89.9\%$ mAP, tăng $0.7$ điểm so với baseline, trong khi
Rank-1 tăng nhẹ từ $95.1\%$ lên $95.2\%$. Mức cải thiện mAP lớn hơn semantic
alignment đơn lẻ, cho thấy việc đánh giá độ tin cậy token cục bộ có ảnh hưởng
rõ đến chất lượng sắp hạng tổng thể. Tuy nhiên, độ trễ tăng từ $4.626$ ms lên
$4.786$ ms, tương đương khoảng $3.45\%$, do reliability head vẫn được sử dụng
khi suy luận.

Việc tăng Rank-1 không nhiều nhưng mAP tăng rõ hơn cho thấy local reliability
chủ yếu cải thiện thứ tự truy hồi ở toàn bộ gallery, thay vì chỉ cải thiện mẫu
đứng đầu. Điều này phù hợp với mục tiêu của reliability branch: giảm đóng góp
của token cục bộ kém tin cậy và ổn định hơn đặc trưng truy hồi trong các trường
hợp có nhiễu hoặc che khuất cục bộ.

\subsection{Hiệu quả của bản kết hợp}

Bản kết hợp semantic alignment và local reliability đạt kết quả tốt nhất:
$90.1\%$ mAP, $95.4\%$ Rank-1, $98.6\%$ Rank-5 và $99.1\%$ Rank-10. So với
baseline, mức cải thiện là $+0.9$ điểm mAP và $+0.3$ điểm Rank-1. Độ trễ trung
bình là $4.783$ ms, gần như tương đương local reliability đơn lẻ và tăng khoảng
$3.39\%$ so với baseline.

Kết quả này cho thấy hai tín hiệu semantic và reliability có tính bổ sung.
Semantic alignment cung cấp ràng buộc part-level trong huấn luyện nhưng không
làm tăng đáng kể chi phí inference. Local reliability bổ sung cơ chế chọn lọc
token đáng tin cậy hơn ở pha truy hồi. Khi kết hợp, mô hình tận dụng được cả
cấu trúc ngữ nghĩa của cơ thể và độ tin cậy cục bộ của patch token, trong khi
vẫn giữ global branch của TransReID làm đường biểu diễn chính.
